% ============================================================
%  5. Architecture Details
% ============================================================
\section{Architecture}
\label{sec:arch}

\subsection{RetVec Embedder}

\RetVec{}~\citep{sander2023retvec} maps any Unicode word $w$ to a
$256$-dimensional vector through a three-layer feed-forward network:
\begin{equation}
  \RetVec(w)
  \;=\;
  \tanh\!\bigl(
    \bW_3\,\text{GELU}(
      \bW_2\,\text{GELU}(
        \bW_1\,\phi(w) + \bb_1
      ) + \bb_2
    ) + \bb_3
  \bigr),
  \label{eq:retvec}
\end{equation}
\noindent where $\phi(w) \in \RR^{384}$ is the fixed character-hash binarisation of word $w$ (16 characters $\times$ 24 bits per character); $\bW_1 \in \RR^{256 \times 384}$, $\bW_2, \bW_3 \in \RR^{256 \times 256}$ are weight matrices; $\bb_1, \bb_2, \bb_3 \in \RR^{256}$ are bias vectors. All weights are pretrained and \textbf{frozen throughout all three stages}. The output lies in $[-1, 1]^{256}$ due to the final $\tanh$ activation.

\subsection{Pre-LayerNorm Transformer Block}

All Transformer layers in both \RetBERT{} and \LightRet{} use the
pre-LayerNorm variant~\citep{xiong2020layer}, which we found more training-stable
during distillation than post-LayerNorm:
\begin{align}
  \bx' &= \bx + \text{MHA}\!\left(\text{LN}(\bx)\right),
  \label{eq:pre_ln_attn}\\
  \bx'' &= \bx' + \text{FFN}\!\left(\text{LN}(\bx')\right),
  \label{eq:pre_ln_ffn}
\end{align}
\noindent where $\bx$ is the block's input vector; $\text{LN}(\cdot)$ is layer normalisation; $\text{MHA}(\cdot)$ is multi-head self-attention (Eq.~\ref{eq:mha}); $\bx'$ is the intermediate output after the attention sub-layer; and $\text{FFN}(\bv) = \bW_2\,\text{GELU}(\bW_1\bv + \bb_1) + \bb_2$ is the position-wise feed-forward network with $\bW_1 \in \RR^{4d \times d}$, $\bW_2 \in \RR^{d \times 4d}$, and $d$ the hidden dimension of the block.

Multi-head attention with $H$ heads and head dimension $d_k = d / H$:
\begin{equation}
  \text{MHA}(\bQ,\bK,\bV)
  \;=\;
  \text{Concat}\!\left(\text{head}_1,\ldots,\text{head}_H\right)\bW^O,
  \label{eq:mha}
\end{equation}
\begin{equation}
  \text{head}_h
  \;=\;
  \text{Attention}\!\left(\bQ\bW^Q_h,\;\bK\bW^K_h,\;\bV\bW^V_h\right)
  \;=\;
  \text{softmax}\!\!\left(\frac{\bQ\bW^Q_h\,(\bK\bW^K_h)^\top}{\sqrt{d_k}}\right)
  \bV\bW^V_h.
  \label{eq:attention}
\end{equation}
\noindent where $\bQ$, $\bK$, $\bV$ are the query, key, and value matrices derived from the input sequence; $\bW^Q_h, \bW^K_h, \bW^V_h \in \RR^{d \times d_k}$ are per-head projection matrices; $d_k = d / H$ is the per-head dimension; $H$ is the total number of attention heads; $\bW^O \in \RR^{d \times d}$ is the output projection applied after concatenation; and $\text{softmax}(\cdot)$ is applied row-wise to produce attention weights.

\subsection{Model Comparison}

Table~\ref{tab:arch} summarises the architectural parameters of
\LightRet{} against its distillation teachers and baselines.

\begin{table}[h]
\centering
\caption{Architecture comparison. \textbf{Vocab} = requires tokenizer vocabulary.}
\label{tab:arch}
\resizebox{\columnwidth}{!}{%
\begin{tabular}{lrrrrll}
\toprule
\textbf{Model} & \textbf{Params} & \textbf{Layers} & \textbf{$d$} & \textbf{$H$} & \textbf{Vocab} & \textbf{Input} \\
\midrule
BiLSTM-CRF       & $\sim$8M   & --  & 256 & --  & Yes & Char+Word  \\
DistilBERT        & 66M        & 6   & 768 & 12  & Yes & Subword    \\
BERT-base         & 110M       & 12  & 768 & 12  & Yes & Subword    \\
CharBERT          & 114M       & 12  & 768 & 12  & Yes & Subword+Char\\
\midrule
\RetBERT{}        & $\sim$86M  & 12  & 768 & 12  & \textbf{No} & \RetVec{} \\
\LightRet{} (ours)& $\sim$4M   & 4   & 256 & 4   & \textbf{No} & \RetVec{} \\
\bottomrule
\end{tabular}
}
\end{table}

\subsection{Three-Stage Architecture Diagram}

Figure~\ref{fig:pipeline} illustrates the full training pipeline.

\begin{figure}[t]
\centering
\begin{tikzpicture}[
  every node/.style={font=\small},
  box/.style={
    draw, rounded corners=4pt, minimum width=2.0cm,
    minimum height=0.6cm, align=center, fill=white, thick
  },
  frozen/.style={box, fill=blue!10},
  train/.style={box,  fill=orange!15},
  loss/.style={box,  fill=green!12, dashed, minimum width=1.5cm},
  arr/.style={-{Latex[length=5pt]}, thick},
  darr/.style={-{Latex[length=5pt]}, thick, dashed, gray}
]

% ─── Stage 1 ───────────────────────────────────
\node[frozen] (bert)    at (0,  4.2) {BERT-base\\(frozen)};
\node[train]  (proj1)   at (0,  3.0) {Linear\\$256{\to}768$};
\node[train]  (rb12)    at (0,  1.8) {$12\times$Transformer\\$d{=}768,H{=}12$};
\node[frozen] (rv1)     at (0,  0.6) {\RetVec{}\\(frozen)};
\node[loss]   (l1)      at (2.2, 3.6) {$\mathcal{L}_1$\\cosine};

\draw[arr] (rv1)  -- (rb12);
\draw[arr] (rb12) -- (proj1);
\draw[darr](proj1) -- node[right,font=\scriptsize]{$\bz^{RB}$} (l1);
\draw[darr](bert)  -- node[right,font=\scriptsize]{$\bz^{B}$}  (l1);

\node[above=0pt of bert, font=\small\bfseries] {Stage 1};

% ─── Stage 2 ───────────────────────────────────
\node[frozen] (rb_t)   at (4.4, 4.2) {\RetBERT{}\\(frozen)};
\node[train]  (proj2)  at (4.4, 3.0) {Linear\\$256{\to}768$};
\node[train]  (tr4)    at (4.4, 1.8) {$4\times$Transformer\\$d{=}256,H{=}4$};
\node[train]  (bigru)  at (4.4, 0.9) {BiGRU\\$128{\times}2$};
\node[frozen] (rv2)    at (4.4, 0.0) {\RetVec{}\\(frozen)};
\node[loss]   (l2)     at (6.6, 3.6) {$\mathcal{L}_2$\\token\\cosine};

\draw[arr] (rv2)   -- (bigru);
\draw[arr] (bigru) -- (tr4);
\draw[arr] (tr4)   -- (proj2);
\draw[darr](proj2) -- node[right,font=\scriptsize]{$\bp$}       (l2);
\draw[darr](rb_t)  -- node[right,font=\scriptsize]{$\bh^{RB}$} (l2);

\node[above=0pt of rb_t, font=\small\bfseries] {Stage 2};

% ─── Stage 3 ───────────────────────────────────
\node[frozen] (lr_t)   at (8.8, 4.2) {\LightRet{}\\teacher (frozen)};
\node[train]  (nerh)   at (8.8, 3.2) {BiLSTM-NER\\$128{\times}2$};
\node[train]  (tr4s)   at (8.8, 1.8) {$4\times$Transformer\\$d{=}256,H{=}4$};
\node[train]  (bigrus) at (8.8, 0.9) {BiGRU\\$128{\times}2$};
\node[frozen] (rv3)    at (8.8, 0.0) {\RetVec{}\\(frozen)};
\node[loss]   (l3)     at (11.0,3.6) {$\mathcal{L}_3$\\class\,+\\distill};

\draw[arr] (rv3)    -- (bigrus);
\draw[arr] (bigrus) -- (tr4s);
\draw[arr] (tr4s)   -- (nerh);
\draw[darr](nerh)   -- (l3);
\draw[darr](lr_t)   -- node[right,font=\scriptsize]{$\bh^T$} (l3);

\node[above=0pt of lr_t, font=\small\bfseries] {Stage 3};

% ─── Noise arrow ───────────────────────────────
\draw[-{Latex}, thick, red!70!black]
  ([xshift=-10pt]rv3.south)
  -- ++(0,-0.45)
  node[below, font=\scriptsize, align=center, red!70!black]
       {noisy\\input $\tilde{\bw}$};

\draw[-{Latex}, thick, blue!70!black]
  ([xshift=10pt]lr_t.south)
  -- ++(0,-3.3)
  node[below, font=\scriptsize, align=center, blue!70!black]
       {clean\\input $\bw$};

\end{tikzpicture}
\caption{The three-stage \LightRet{} training pipeline.
  Blue boxes are frozen; orange boxes are trainable.
  In Stage~3 the teacher (left) receives clean text $\bw$,
  while the student (right) receives noisy text $\tilde{\bw}$.}
\label{fig:pipeline}
\end{figure}

\subsection{Parameter Count Breakdown}

For \LightRet{} the trainable parameters are distributed as follows:

\begin{table}[h]
\centering
\caption{\LightRet{} parameter breakdown.}
\label{tab:params}
\begin{tabular}{lr}
\toprule
\textbf{Component} & \textbf{Parameters} \\
\midrule
\RetVec{} (frozen, not counted) & 256-d output \\
BiGRU ($2 \times 3 \times 256 \times 128 \times 2$) & $\sim$787K \\
$4\times$ Transformer ($d{=}256, H{=}4, \text{FFN}{=}1024$) & $\sim$2.6M \\
BiLSTM NER head ($d{=}256{\to}256$) & $\sim$526K \\
Linear classifier ($256{\to}9$) & $\sim$2.3K \\
\midrule
\textbf{Total} & $\bm{\sim}$\textbf{3.9M} \\
\bottomrule
\end{tabular}
\end{table}
